\chapter{Combined CSP \& DSP approach to categorising signal families}

\section{Cyclostationary analysis}

A signal with the property $x(t) = x(t + T)\ \forall\, t,\ T>0$ is called \textit{strictly periodic} (or \textit{periodic pointwise}).
A composition of periodic signals $h(t) = \sum_{i=0}^N x_i(t)$ is also periodic $\iff \tfrac{T_i}{T_j} \in \mathbb{Q}\ \forall\, i,j$ (i.e., the periods are commensurable).
The period of $h(t)$ can therefore be described as
\[
T = \mathrm{lcm}(T_0, T_1, \dots, T_N).
\]

A \textit{wide-sense cyclostationary} (WSC) signal is defined by
\[
\mathbb{E}[x(t)] = \mathbb{E}[x(t+T)], \quad R_x(t,\tau) = R_x(t+T,\tau),
\]
where $R_x(t,\tau) = \mathbb{E}[x(t)x^*(t+\tau)]$.
That is, both the mean $\bar{x}(t)$ and the autocorrelation function $R_x(t,\tau)$ are periodic in $t$.
More generally, a signal is cyclostationary \textit{iff} there exists some $n$th-order nonlinear transformation of $x(t)$ that produces a finite additive sinusoidal component~\cite{}.

The definition of the \textit{spectral parameter}~\cite{} is
\[
\hat{x}(\alpha) \;=\; \lim_{T \to \infty} \frac{1}{T}
\int_{-T/2}^{T/2} x(t)\,e^{-j2\pi \alpha t}\,dt.
\]

For $x(t) = a\cos(2\pi \beta t + \theta)$:
\[
x(t) = \tfrac{a}{2}\big(e^{j(2\pi \beta t + \theta)} + e^{-j(2\pi \beta t + \theta)}\big).
\]

Substituting into the definition:
\[
\hat{x}(\alpha) = \lim_{T \to \infty} \frac{a}{2T}
\left[
e^{j\theta}\int_{-T/2}^{T/2} e^{j2\pi(\beta-\alpha)t}\,dt \;+\;
e^{-j\theta}\int_{-T/2}^{T/2} e^{-j2\pi(\beta+\alpha)t}\,dt
\right].
\]

Recalling that
\[
\int_{-T/2}^{T/2} e^{j2\pi f t}\,dt = \frac{\sin(\pi f T)}{\pi f},
\]
we obtain
\[
\hat{x}(\alpha) = \frac{a}{2} \lim_{T \to \infty} \frac{1}{T}\Bigg[
e^{j\theta}\frac{\sin(\pi(\beta-\alpha)T)}{\pi(\beta-\alpha)}
+ e^{-j\theta}\frac{\sin(\pi(\beta+\alpha)T)}{\pi(\beta+\alpha)}
\Bigg].
\]

Equivalently,
\[
\hat{x}(\alpha) = \frac{a}{2}\lim_{T \to \infty}\left[e^{j\theta}\,\mathrm{sinc}((\beta-\alpha)T)
+ e^{-j\theta}\,\mathrm{sinc}((\beta+\alpha)T)\right],
\]
where $\mathrm{sinc}(x) = \frac{\sin(\pi x)}{\pi x}$.

If $\alpha = \beta$:
\[
\hat{x}(\alpha) = \tfrac{a}{2}e^{j\theta}.
\]

If $\alpha = -\beta$:
\[
\hat{x}(\alpha) = \tfrac{a}{2}e^{-j\theta}.
\]

Otherwise:
\[
\hat{x}(\alpha) = 0.
\]

Thus:
\[
\hat{x}(\alpha) =
\begin{cases}
	\dfrac{a}{2}e^{j\theta}, & \alpha = \beta,\\[6pt]
	\dfrac{a}{2}e^{-j\theta}, & \alpha = -\beta,\\[6pt]
	0, & \text{otherwise}.
\end{cases}
\]

This demonstrates that the time-averaged inner product survives only when the test frequency $\alpha$ matches a sinusoidal component $\beta$ of the signal.

\subsection{Relation between Strict Periodicity and Cyclostationarity}

Every strictly periodic signal is trivially wide-sense cyclostationary.
Indeed, if $x(t) = x(t+T)$, then
\[
\mathbb{E}[x(t)] = \mathbb{E}[x(t+T)], \quad
R_x(t,\tau) = \mathbb{E}[x(t)x^*(t+\tau)] = \mathbb{E}[x(t+T)x^*(t+T+\tau)] = R_x(t+T,\tau).
\]
Thus:
\[
x(t)\ \text{periodic} \;\implies\; x(t)\ \text{WSC (cyclostationary)}.
\]

The converse does not hold: cyclostationary signals need not be strictly periodic pointwise.
For example, an amplitude-modulated process $x(t) = m(t)\cos(2\pi f_c t)$ with random data $m(t)$ can be cyclostationary with period $1/f_c$, while not being pointwise periodic.

\section{Novel combined CSP, DSP approach to creating FPGA cache heirarchies}

\[
\text{A function } f: A \to B \text{ is \emph{foldable} iff }
\exists \, A_0 \subseteq A, \, T = \{ \tau_i : A \to A \}_{i \in I}
\]
such that
\[
\forall x \in A, \, \exists \, \tau \in T \text{ with } \tau(x) \in A_0
\quad \text{and} \quad f(x) = f \circ \tau(x).
\]\\

Where $A_0$ is called the \textit{fundamental domain} \& $T$, the set of transformations mapping from $A_0$ to $A$, the \textit{recovery method}.\\

Motivating this definition with a classic example, consider $\sin(x)$ which has period $2\pi$ and can be represented with fundamental domain $\frac{\pi}{4}$ where the recovery method

\[
T =
\begin{cases}
	x, & 0 \leq x \leq \frac{\pi}{4},\\[6pt]
	\pi - x, & \frac{\pi}{4} < x \leq \frac{\pi}{2},\\[6pt]
	x - \pi, & \frac{\pi}{2} < x \leq \frac{3\pi}{2},\\[6pt]
	2\pi - x, & \frac{3\pi}{2} < x \leq 2\pi
\end{cases}
\]\\

\[
f(x) =
\begin{cases}
	f \circ \tau(x), & 0 \leq x \leq \pi,\\[6pt]
	-f \circ \tau(x), & \pi < x \leq 2\pi
\end{cases}
\]

Can be used to store effectively $4\times$ less memory $M$ than the full period \textbf{OR} with the same memory footprint but $4\times$ greater granularity $G$. In general, the ratio between $M$ \& $G$, $\frac{M}{G}$, is called the \textit{precision-memory gain tradeoff}. $M = \frac{P}{|A_0|}$, $P$ is the period of the function.\\

\textbf{Note:} it is important that the recovery method $f \circ \tau(x)$ is relatively efficient / portable to an FPGA \& that the tradeoff between the computational complexity of the recovery method \& the memory gain $M$ is not necessarily linear E.g. $\tau(x) \not\propto M$. A relevant example explored later on \cite{} is the function $\arccos(x)$ which requires the computation of $\sqrt{x + 1}$ on an FPGA in order to do the lookup on the reduced domain $\arccos: [0, \frac{1}{\sqrt{2}}] \to [-1, 1]$ \\

\subsection*{Group-Theoretic Interpretation}
The set of transformations $T$ can be understood as the action of a \emph{symmetry group} $G$ on the domain $A$. In the sine example, the relevant group is the \emph{dihedral group} $D_4$, which consists of rotations and reflections of a square \cite{armstrong1988groups, gallian2021contemporary}.
Each $\tau \in T$ corresponds to an element of $D_4$ acting on the domain of $\sin(x)$, with reflections accounting for sign changes and rotations corresponding to translations by $\pi/2$.

Thus, the concept of foldability is equivalent to requiring that the function $f$ is invariant under the action of a group $G$, up to sign or phase factors. The \emph{fundamental domain} $A_0$ is then a representative region of $A$ under this group action, and folding reduces storage by exploiting orbit representatives of $G$.

In group-theoretic language, the prior example is invariant under the group action up to a sign factor, i.e.,
\[
\sin(x) = \epsilon(\tau) \, \sin(\tau(x)), \quad \epsilon(\tau) \in \{+1, -1\}, \ \tau \in G.
\]
The fundamental domain $A_0$ corresponds to a set of orbit representatives of $G$ acting on $A$, and the recovery method reconstructs the full signal from this reduced domain.

\subsection*{Link to Cyclostationary Signals}

Foldable functions are naturally related to cyclostationarity. A cyclostationary signal $x(t)$ has statistics that are periodic with period $T$.\\

For a foldable function like $\sin(x)$, the fundamental domain $A_0$ can be interpreted as a segment over which the first- and second-order statistics repeat under the action of $G$. In other words, folding exploits the same periodic structure that cyclostationary analysis detects: the mean and autocorrelation within $A_0$ are sufficient to reconstruct the signal statistics over the entire domain.

Thus, foldability can be seen as a deterministic, structural analogue of cyclostationarity: instead of averaging over time to detect periodic statistics, we use group-theoretic transformations to reconstruct the full signal from a reduced representative domain.

Let \( f: A \to B \) be a foldable function with fundamental domain \( A_0 \subseteq A \) and recovery transformations \( T = \{ \tau_i : A \to A \}_{i \in I} \) such that
\[
\forall x \in A, \exists \tau \in T \text{ with } \tau(x) \in A_0 \quad \text{and} \quad f(x) = f \circ \tau(x).
\]
This implies that the function \( f \) is invariant under the group action of \( G \) generated by \( T \), and the entire domain \( A \) can be reconstructed from \( A_0 \) using the transformations in \( T \).

Now, consider the autocorrelation function of \( f \):
\[
R_f(x, \Delta) = \mathbb{E}[f(x) f^*(x + \Delta)].
\]
Since \( f(x) = f(\tau(x)) \) for some \( \tau \in T \), we have
\[
R_f(x, \Delta) = \mathbb{E}[f(\tau(x)) f^*(\tau(x + \Delta))].
\]
Because \( \tau(x), \tau(x + \Delta) \in A_0 \) (up to another element of the group), all evaluations of \( R_f(x, \Delta) \) are determined entirely by the behavior on \( A_0 \). Therefore,
\[
R_f(x, \Delta) = R_f(\tau(x), \Delta), \quad \forall x \in A, \ \tau \in T, \ \tau(x) \in A_0.
\]
Thus, the fundamental domain \( A_0 \) captures all statistical information of the function \( f \), analogous to how the fundamental domain in cyclostationary signal processing captures the periodic statistical properties of a signal.

\subsection{Procedure to storing arbitrary signals}

A procedure to \textit{caching} an arbitrary signal $x(t)$ is the following:

\begin{enumerate}
    \item Decompose the signal into a series of harmonics (for analysis purposes).
          In practice, this can be approximated via autocorrelation or spectral analysis rather than explicit FFT computation.
    \item Use autocorrelation (ACF) analysis on the fundamental harmonic to estimate the period of a noisy signal within a specified SNR tolerance.
    \item Iteratively refine the period estimate using heuristics, e.g., minimize spectral leakage if the period is undershot, or adjust lags using peak prominence methods.
    \item Determine if the signal is \emph{foldable} by checking whether there exists a fundamental domain $A_0$ and a set of transformations $T$ such that
    \[
        \forall x \in A, \ \exists \tau \in T \text{ with } \tau(x) \in A_0 \text{ and } f(x) = f \circ \tau(x).
    \]
    This step is limited to signals with known symmetries (e.g., trigonometric functions, piecewise monotonic functions).
    \item Construct a recovery method that maps any $x \in A$ to $A_0$ via the transformations in $T$, taking into account the desired precision-memory gain tradeoff based on control-theoretic or resource-demand considerations.
    \item Implement the recovery method in a synthesizable and efficient form for FPGA deployment, ensuring that the computational complexity of applying $\tau(x)$ does not outweigh the memory savings.
\end{enumerate}

Methods for \textit{(1), (2) and (3)}, previously discussed and implemented, are the subject of extensive optimization effort. Steps \textit{(4), (5), and (6)}, however, are considerably more difficult to achieve in practice on a completely unrestricted problem domain. Indeed, this paper focuses on a subset of signal definitions that can be reliably categorized, as well as on certain elementary signals that are relatively commonplace and often appear as components within more stochastic processes.

\section{Synthesizable construction of the recovery method, $T$, for FPGA's}

\subsection*{Efficient FPGA Implementation of \texorpdfstring{$\arccos(x)$}{arccos(x)} Using Quarter-Domain Folding}

In previous sections, we discussed foldable functions and their use in memory-efficient lookup. In this section, we consider the case of $\arccos(x)$, which initially seems trivial but actually presents valuable insights for efficient computation on hardware.

\subsubsection*{Fundamental Domain Recovery}

The key identity used for recovery from a quarter-domain representation is:
\begin{align}
\tau(x) &= \frac{1}{\sqrt{2}}\sqrt{x + 1}, \nonumber \\
f(x) &= 2 \cdot \arccos(\tau(x)) \\
\implies \arccos(x) &= 2 \, \arccos\!\left(\frac{1}{\sqrt{2}} \, \sqrt{x+1}\right).
\end{align}

Applying a third-order Taylor series expansion to $\sqrt{x+1}$ yields:
\[
16 \sqrt{x+1} \approx x^3 - 2 x^2 + 8 x + 16.
\]

Unlike the $\arcsin$ third-order polynomial, this polynomial can be factored over a finite field:
\[
(x+1)(x^2+x+1).
\]

To recover the original polynomial, we add the remaining series:
\[
T_2 = -4 x^2 + 7x + 15.
\]

\subsubsection*{Efficient Hardware Computation}

Computing the first term:
\[
T_1 = (x+1)(x^2+x+1)
\]

(\textbf{Note} that the second term $T_2$ can be computed using precomputed values)
\[
T_2 = -4 x^2 + 7 x + 15 = -4x^2 + 8x + 16 -(x + 1) = -(x^2 << 2) + ((x+1) + 1) << 3) - (x+1)
\]

Which can be implemented as follows:

\begin{verbatim}
// Begin computation of tau(x) for arccos(x)
// arccos in [0, 1/sqrt(2)], Q2.21 format (24 bit data-width)

// Clock cycle 1: compute shared terms
t1 = x + 1;
x_squared = x * x;

// Clock cycle 2: compute T1
t1 = t1 * (x_squared + t1);

// Clock cycle 2: concurrent with T1, compute T2
t2 = -(x_squared << 2) + ((t1 + 1) << 3) - t1;

// Clock cycle 3: add finite-field-factorization correction term, T2, to T1
tau = t1 + t2;

// Clock cycle 4: right-shift by 4 to get sqrt(x + 1), multiply by 1/sqrt(2) term
localparam inv_sqrt2_q = ...; // 1/sqrt(2) stored in Q2.21 format
tau = (tau >> 4) * inv_sqrt2_q;
\end{verbatim}

\begin{verbatim}
// Do the other required computations to get T

// Clock cycle 5/6: quantization of tau(x) to indice present in BRAM, LUT lookup
logic [23:0] T;
acos_tau_quantize f (.tau(tau), .f(T), .clk(clk));

// Clock cycle 7: final result
T = T << 1;

\end{verbatim}

\subsubsection*{Resource Analysis and Benefits}

The total cost for this third-order implementation of $\tau(x)$:
\begin{itemize}
    \item 3 multiplications (1 for shared term, 1 for $T_1$, 1 for multiplication by $\frac{1}{\sqrt{2}}$)
    \item 6 additions/subtractions (1 for shared term, 1 for $T_1$, 3 for $T_2$, 1 for the final sum)
    \item 4 bit-shifts and 1 bitwise inversion (sign flip)
\end{itemize}

\paragraph{Advantages:}
\begin{enumerate}
    \item No transcendental operations in the argument apart from the $\frac{1}{\sqrt{2}}$ multiplier, simplifying quantization.
    \subitem quantization step can be bypassed to save a clock cycle by computing the $\arccos$ lut values along a non-linear axis I.e. according to the distribution $tau(x)$
    \item Only a single DSP unit is required for all multiplies if the data width allows single-cycle multiplication.
    \item No additional storage for constants required except for $\frac{1}{\sqrt{2}}$.
    \item The computation is easily pipelined.
    \item Bit-growth is predictable, scaling relative to $\deg(\tau(x))$
\end{enumerate}
